\documentclass[11pt,a4paper,UTF8]{ctexart}
\usepackage{geometry}
\geometry{left=18mm,right=18mm,top=24mm,bottom=24mm,headheight=22pt,footskip=12mm}
\usepackage{amsmath,amssymb,mathtools,bm,esint,extarrows,centernot}
\usepackage{xcolor}
\usepackage{tcolorbox}
\tcbuselibrary{skins,breakable}
\usepackage{fancyhdr}
\usepackage{titlesec}
\usepackage{hyperref}
\usepackage{tikz}
\usepackage{booktabs}
\usepackage{tabularx}

% --- Color Palette (Purple & DeepPink Academic Theme) ---
\definecolor{DeepPurple}{HTML}{4C1D95}
\definecolor{TitlePurple}{HTML}{581C87}
\definecolor{SubPurple}{HTML}{6B21A8}
\definecolor{DeepPink}{HTML}{FF1493}
\definecolor{LightPinkBg}{HTML}{FFF5F8}
\definecolor{LightPurpleBg}{HTML}{F8F5FF}
\definecolor{GoldAccent}{HTML}{D97706}
\definecolor{DarkText}{HTML}{1F2937}

\hypersetup{
    colorlinks=true,
    linkcolor=TitlePurple,
    citecolor=DeepPink,
    urlcolor=SubPurple,
    pdfborder={0 0 0}
}

% --- Typography & Spacing ---
\linespread{1.3}
\setlength{\parskip}{0.35em plus 0.1em minus 0.1em}
\setlength{\parindent}{0pt}

% --- Section Titles Styling ---
\titleformat{\section}
  {\Large\bfseries\color{TitlePurple}}{\thesection}{1em}{}
  [\vspace{1mm}{\color{TitlePurple!30}\hrule height 1pt}]
\titleformat{\subsection}
  {\large\bfseries\color{SubPurple}}{\thesubsection}{0.8em}{}
\titleformat{\subsubsection}
  {\normalsize\bfseries\color{DeepPurple}}{\thesubsubsection}{0.6em}{}

% --- tcolorbox Environments ---
% 1. Theorem / Formula / Method / Analysis Box (Deep Purple)
\newtcolorbox{thmbox}[1][]{
    enhanced,
    breakable,
    colback=LightPurpleBg,
    colframe=TitlePurple,
    arc=3pt,
    boxrule=0pt,
    leftrule=3.5pt,
    left=10pt,
    right=10pt,
    top=8pt,
    bottom=8pt,
    title=#1,
    coltitle=white,
    colbacktitle=TitlePurple,
    fonttitle=\bfseries\small,
    attach boxed title to top left={yshift=-2mm, xshift=4mm},
    boxed title style={boxrule=0pt, arc=2pt}
}

% 2. Solution / Formula / Derivation Box (DeepPink)
\newtcolorbox{solbox}[1][]{
    enhanced,
    breakable,
    colback=LightPinkBg,
    colframe=DeepPink,
    arc=3pt,
    boxrule=0pt,
    leftrule=3.5pt,
    left=10pt,
    right=10pt,
    top=8pt,
    bottom=8pt,
    title=#1,
    coltitle=white,
    colbacktitle=DeepPink,
    fonttitle=\bfseries\small,
    attach boxed title to top left={yshift=-2mm, xshift=4mm},
    boxed title style={boxrule=0pt, arc=2pt}
}

% 3. Learning Goals / Summary Box
\newtcolorbox{targetbox}[1][]{
    enhanced,
    breakable,
    colback=LightPurpleBg,
    colframe=DeepPurple,
    arc=3pt,
    boxrule=1pt,
    left=10pt,
    right=10pt,
    top=8pt,
    bottom=8pt,
    title=#1,
    coltitle=white,
    colbacktitle=DeepPurple,
    fonttitle=\bfseries\small,
    attach boxed title to top left={yshift=-2mm, xshift=4mm},
    boxed title style={boxrule=0pt, arc=2pt}
}

\begin{document}

% ------------------- 讲义封面 (Cover Page) -------------------
\begin{titlepage}
\thispagestyle{empty}
\begin{tikzpicture}[remember picture,overlay]
    \fill[DeepPurple] (current page.north west) rectangle ([yshift=-7cm]current page.north east);
    \draw[DeepPink, line width=3.5pt] ([yshift=-7cm]current page.north west) -- ([yshift=-7cm]current page.north east);
    \draw[GoldAccent, line width=1pt] ([yshift=-7.12cm]current page.north west) -- ([yshift=-7.12cm]current page.north east);
    
    \node[anchor=north east, opacity=0.12, text=white] at ([xshift=-1.5cm, yshift=-0.5cm]current page.north east) {
        \fontsize{70}{70}\selectfont\bfseries 2027
    };
\end{tikzpicture}

\vspace*{0.8cm}
\begin{center}
    {\color{white}\fontsize{26}{32}\selectfont\bfseries 2027 考研数学(二) 核心讲义}\\[0.6cm]
    {\color{white}\fontsize{13}{18}\selectfont\bfseries 全国硕士研究生招生考试 · 统考数学(二) 核心精讲全书}\\[1.5cm]
\end{center}

\vspace{3.5cm}
\begin{center}
    {\color{GoldAccent}\large\bfseries $\blacktriangleright$ 基础强化 · 核心题解 · 考点深水区 $\blacktriangleleft$}\\[0.8cm]
    {\color{TitlePurple}\fontsize{24}{28}\selectfont\bfseries 第八章 一元函数积分学的概念与性质}\\[0.6cm]
    {\color{DeepPink}\large\bfseries —— 课后精选练习题与全过程推导 ——}\\[1.5cm]
\end{center}

\vfill

\begin{center}
\begin{tcolorbox}[
    colback=white,
    colframe=DeepPurple,
    width=0.88\textwidth,
    arc=4pt,
    boxrule=1.2pt,
    halign=center,
    drop shadow=gray!30
]
    \vspace{3mm}
    {\bfseries\large 编著团队：EscoffierZhou}\\[2.5mm]
    {\bfseries\color{TitlePurple} 目标院校：中国科学院大学 (UCAS) 计算机专硕 (22408)}\\[2.5mm]
    {\small\color{gray} 备考铁律：手算真题数字 · 错题白纸重做 · 408 客观题为王}\\[2.5mm]
    {\small 编制版本：2027 届首发版 · \today}
    \vspace{2mm}
\end{tcolorbox}
\end{center}
\vspace{0.8cm}
\end{titlepage}

% ------------------- 目录与页面主体配置 -------------------
\pagestyle{fancy}
\fancyhf{}
\fancyhead[L]{\small\color{gray}2027 考研数学(二) · 核心讲义系列}
\fancyhead[R]{\small\color{TitlePurple}\nouppercase{\leftmark}}
\fancyfoot[C]{\small\color{gray}— 第 \thepage\ 页 —}
\fancyfoot[R]{\small\color{gray}UCAS 22408}
\renewcommand{\headrulewidth}{0.5pt}

\pagenumbering{Roman}
\tableofcontents
\newpage
\pagenumbering{arabic}


\section{8.Homework}


\subsection{习题 1: 三角不等式与定积分保号性比较大小}

设 $I_1 = \int_0^{\frac{\pi}{2}} \frac{\sin x}{x}\,\mathrm{d}x$，$I_2 = \int_0^{\frac{\pi}{2}} \frac{x}{\sin x}\,\mathrm{d}x$，则 $(\quad)$

(A) $I_1 > I_2 > 1 \qquad$ (B) $1 > I_1 > I_2 \qquad$ (C) $I_2 > I_1 > 1 \qquad$ (D) $1 > I_2 > I_1$

\textbf{\color{DeepPurple}【思路点拨】} 在开区间 $(0, \frac{\pi}{2})$ 上利用经典几何三角不等式 $\sin x < x$ 比较被积函数大小，结合若尔当不等式 $\frac{\sin x}{x} > \frac{2}{\pi}$ 对 $I_1$ 进行下界估计


\begin{solbox}[分步推导与答题规范]
\textbf{[1]} 比较 $I_1$ 与 $I_2$ 的大小：

当 $x \in (0, \frac{\pi}{2})$ 时，经典不等式恒成立：$\sin x < x$。

两边同除以正数得：


\begin{equation*}
\frac{x}{\sin x} > 1 > \frac{\sin x}{x} \implies \frac{x}{\sin x} > \frac{\sin x}{x}
\end{equation*}


由定积分保号性可知：


\begin{equation*}
I_2 = \int_0^{\frac{\pi}{2}} \frac{x}{\sin x}\,\mathrm{d}x > \int_0^{\frac{\pi}{2}} \frac{\sin x}{x}\,\mathrm{d}x = I_1
\end{equation*}


由此排除选项 (A) 与 (D)。

\textbf{[2]} 对 $I_1$ 进行若尔当不等式下界放缩：

在 $(0, \frac{\pi}{2})$ 上，根据若尔当(Jordan)不等式 $\frac{\sin x}{x} > \frac{2}{\pi}$：


\begin{equation*}
I_1 = \int_0^{\frac{\pi}{2}} \frac{\sin x}{x}\,\mathrm{d}x > \int_0^{\frac{\pi}{2}} \frac{2}{\pi}\,\mathrm{d}x = \frac{2}{\pi} \cdot \frac{\pi}{2} = 1
\end{equation*}


\textbf{[3]} 综合得出结论：

综合两步可得：$I_2 > I_1 > 1$。

\textbf{[4]} 故正确选项为 \textbf{(C)}。
\end{solbox}


\subsection{习题 2: 分段函数原函数存在性与变限积分可导性判定}

设 $f(x) = \begin{cases} x, & x < 0 \\ \frac{1}{2}, & x = 0 \\ 1, & x > 0 \end{cases}$，则下列命题：

\textbf{(1)}  $f(x)$ 在 $[-1, 1]$ 上有原函数；
\textbf{(2)}  $f(x)$ 在 $[-1, 1]$ 上可积；
\textbf{(3)}  $F(x) = \int_0^x f(t)\,\mathrm{d}t$ 在 $x=0$ 处可导；
\textbf{(4)}  $F(x) = \int_0^x f(t)\,\mathrm{d}t$ 在 $x=0$ 处连续但不可导。

正确命题的个数为 $(\quad)$

(A) 1 \qquad (B) 2 \qquad (C) 3 \qquad (D) 4

\textbf{\color{DeepPurple}【思路点拨】} 检查分界点左右极限识别间断点类型，运用达布介值定理和变上限积分平滑性定理逐项分析


\begin{solbox}[分步推导与答题规范]
\textbf{[1]} 考察间断点类型：

$\lim\limits_{x \to 0^-} f(x) = 0$，$\lim\limits_{x \to 0^+} f(x) = 1$。

由于左右极限均存在但不相等，故 $x = 0$ 是 $f(x)$ 的第一类跳跃间断点。

\textbf{[2]} 分析命题(1)与(2)：

命题(1)错误: 含有第一类间断点的函数在包含该点的区间上必无原函数(达布定理)；

命题(2)正确: $f(x)$ 在 $[-1, 1]$ 上有界且仅有 1 个跳跃间断点，由充分条件知定积分必定存在(即可积)。

\textbf{[3]} 分析命题(3)与(4)：

变上限积分 $F(x)$ 必处处连续；

在跳跃间断点处，左右导数分别等于被积函数的左右极限：


\begin{equation*}
F'_-(0) = f(0^-) = 0, \quad F'_+(0) = f(0^+) = 1
\end{equation*}


因为 $F'_-(0) \neq F'_+(0)$，故 $F(x)$ 在 $x = 0$ 处不可导。

故命题(3)错误，命题(4)正确。

\textbf{[4]} 综合得出正确命题数量：

正确的命题为 (2) 和 (4)，共 2 个。

\textbf{[5]} 故正确选项为 \textbf{(B)}。
\end{solbox}


\subsection{习题 3: 含对数因子的瑕积分与无穷反常积分双奇点审敛}

若反常积分 $\int_0^{+\infty} \frac{\ln x}{(1+x)x^{1-p}}\,\mathrm{d}x$ 收敛，则 $(\quad)$

(A) $p < 1 \qquad$ (B) $p > 1 \qquad$ (C) $0 < p < 1 \qquad$ (D) $0 \leqslant p < 1$

\textbf{\color{DeepPurple}【思路点拨】} 识别积分的两个奇点 $x = 0$(瑕点)与 $x = +\infty$(无穷)，拆分为两段分别由低次项和高次项主导，应用对数阶数压制基准判别


\begin{solbox}[分步推导与答题规范]
\textbf{[1]} 拆分双奇点：


\begin{equation*}
I = \int_0^1 \frac{\ln x}{(1+x)x^{1-p}}\,\mathrm{d}x + \int_1^{+\infty} \frac{\ln x}{(1+x)x^{1-p}}\,\mathrm{d}x = I_1 + I_2
\end{equation*}


原积分收敛当且仅当 $I_1$ 与 $I_2$ 同时收敛。

\textbf{[2]} 分析瑕积分 $I_1$ ($x \to 0^+$)：

当 $x \to 0^+$ 时，$1+x \to 1$，被积函数主部为：


\begin{equation*}
\frac{\ln x}{(1+x)x^{1-p}} \sim \frac{\ln x}{x^{1-p}}
\end{equation*}


由瑕积分对数基准型可知，$\int_0^1 \frac{\vert{}\ln x\vert{}}{x^k}\,\mathrm{d}x$ 收敛的充要条件为幂次 $k < 1$。

这里 $k = 1 - p$，因此要求：


\begin{equation*}
1 - p < 1 \implies p > 0
\end{equation*}


\textbf{[3]} 分析无穷区间积分 $I_2$ ($x \to +\infty$)：

当 $x \to +\infty$ 时，$1+x \sim x$，被积函数主部为：


\begin{equation*}
\frac{\ln x}{(1+x)x^{1-p}} \sim \frac{\ln x}{x \cdot x^{1-p}} = \frac{\ln x}{x^{2-p}}
\end{equation*}


由无穷区间对数基准型可知，$\int_1^{+\infty} \frac{\ln x}{x^k}\,\mathrm{d}x$ 收敛的充要条件为幂次 $k > 1$。

这里 $k = 2 - p$，因此要求：


\begin{equation*}
2 - p > 1 \implies p < 1
\end{equation*}


\textbf{[4]} 求交集得出参数范围：

两端必须同时收敛，取交集得：$0 < p < 1$。

\textbf{[5]} 故正确选项为 \textbf{(C)}。
\end{solbox}


\subsection{习题 4: 定积分定义求解连和式极限}

$\lim\limits_{n \to \infty} \left( \frac{1}{\sqrt{n^2+n}} + \frac{1}{\sqrt{n^2+2n}} + \frac{1}{\sqrt{n^2+3n}} + \dots + \frac{1}{\sqrt{n^2+n^2}} \right) = \underline{\hspace{2.5cm}}$。

\textbf{\color{DeepPurple}【思路点拨】} 通项中分母提取公因式 $n$，分离出 $\frac{1}{n}$ 凑自变量 $\frac{i}{n}$，严格应用定积分定义转化为连续定积分求解


\begin{solbox}[分步推导与答题规范]
\textbf{[1]} 求和通项分析与提取 $\frac{1}{n}$：

将求和式用 $\Sigma$ 紧凑表示：


\begin{equation*}
\text{原式} = \lim_{n \to \infty} \sum_{i=1}^n \frac{1}{\sqrt{n^2+ni}}
\end{equation*}


分母提取公因式 $n$：


\begin{equation*}
= \lim_{n \to \infty} \sum_{i=1}^n \frac{1}{n\sqrt{1 + \frac{i}{n}}} = \lim_{n \to \infty} \frac{1}{n}\sum_{i=1}^n \frac{1}{\sqrt{1 + \frac{i}{n}}}
\end{equation*}


\textbf{[2]} 定积分定义转化：

令 $\frac{i}{n} \to x$，$\frac{1}{n} \to \mathrm{d}x$，求和指标 $i=1 \sim n$ 对应积分区间 $[0, 1]$：


\begin{equation*}
= \int_0^1 \frac{1}{\sqrt{1+x}}\,\mathrm{d}x
\end{equation*}


\textbf{[3]} 计算定积分值：


\begin{equation*}
= \left. 2\sqrt{1+x} \right\vert{}_0^1 = 2\sqrt{2} - 2
\end{equation*}


\textbf{[4]} 故最终答案为 $2\sqrt{2} - 2$。
\end{solbox}


\subsection{习题 5: 夹逼准则化归定积分定义求解微扰和式极限}

$\lim\limits_{n \to \infty} \sum\limits_{i=1}^n \frac{\sin \frac{i\pi}{n}}{n + \frac{1}{i}} = \underline{\hspace{2.5cm}}$。

\textbf{\color{DeepPurple}【思路点拨】} 分母微扰项 $\frac{1}{i}$ 破坏了标准 $\frac{i}{n}$ 结构，利用有界放缩将其夹在两个标准定积分定义式之间，结合夹逼准则求出极限


\begin{solbox}[分步推导与答题规范]
\textbf{[1]} 建立不等式夹逼：

由于对所有 $1 \leqslant i \leqslant n$，恒有 $0 < \frac{1}{i} \leqslant 1$，且 $\sin\frac{i\pi}{n} \geqslant 0$，故有不等式：


\begin{equation*}
\frac{\sin \frac{i\pi}{n}}{n+1} \leqslant \frac{\sin \frac{i\pi}{n}}{n + \frac{1}{i}} \leqslant \frac{\sin \frac{i\pi}{n}}{n}
\end{equation*}


\textbf{[2]} 对不等式求和并计算两端极限：

右端极限：


\begin{equation*}
\lim_{n \to \infty} \frac{1}{n}\sum_{i=1}^n \sin\left(\frac{i}{n}\pi\right) = \int_0^1 \sin(\pi x)\,\mathrm{d}x = \left. \left( -\frac{1}{\pi}\cos(\pi x) \right) \right\vert{}_0^1 = -\frac{1}{\pi}(-1 - 1) = \frac{2}{\pi}
\end{equation*}


左端极限：


\begin{equation*}
\lim_{n \to \infty} \sum_{i=1}^n \frac{\sin \frac{i\pi}{n}}{n+1} = \lim_{n \to \infty} \frac{n}{n+1} \cdot \left[ \frac{1}{n}\sum_{i=1}^n \sin\left(\frac{i}{n}\pi\right) \right] = 1 \cdot \frac{2}{\pi} = \frac{2}{\pi}
\end{equation*}


\textbf{[3]} 应用夹逼准则得出结论：

由夹逼准则，原极限等于 $\frac{2}{\pi}$。

\textbf{[4]} 故最终答案为 $\frac{2}{\pi}$。
\end{solbox}


\subsection{习题 6: 周期锯齿波变上限积分的单侧导数计算}

设 $F(x) = \int_0^x (t - [t])\,\mathrm{d}t$，其中 $[x]$ 表示不超过 $x$ 的最大整数，则 $F'_-(1) + F'_+(1) = \underline{\hspace{2.5cm}}$。

\textbf{\color{DeepPurple}【思路点拨】} 被积函数为周期锯齿波，整数点均为跳跃间断点；变上限积分处处连续，分界点处的单侧导数直接对应被积函数的单侧极限


\begin{solbox}[分步推导与答题规范]
\textbf{[1]} 分析被积函数的间断点：

记被积函数为 $f(t) = t - [t]$(表示 $t$ 的小数部分，周期为 1 的锯齿波)。

在整数分界点处，$f(t)$ 产生第一类跳跃间断点，其变上限积分 $F(x)$ 处处连续。

\textbf{[2]} 计算分界点 $x = 1$ 处的单侧导数：

变上限积分的单侧导数等于被积函数的单侧极限：

左导数：


\begin{equation*}
F'_-(1) = \lim_{t \to 1^-} f(t) = \lim_{t \to 1^-} (t - [t]) = \lim_{t \to 1^-} (t - 0) = 1
\end{equation*}


右导数：


\begin{equation*}
F'_+(1) = \lim_{t \to 1^+} f(t) = \lim_{t \to 1^+} (t - [t]) = \lim_{t \to 1^+} (t - 1) = 0
\end{equation*}


\textbf{[3]} 两单侧导数求和：


\begin{equation*}
F'_-(1) + F'_+(1) = 1 + 0 = 1
\end{equation*}


\textbf{[4]} 故最终答案为 $1$。
\end{solbox}


\subsection{习题 7: 含对数幂次复合反常积分的敛散性讨论}

讨论反常积分 $\int_2^{+\infty} \frac{1}{x \ln^p x}\,\mathrm{d}x$ 的敛散性，其中 $p$ 为任意实数。

\textbf{\color{DeepPurple}【思路点拨】} 利用凑微分法将对数因子换元为基本变量 $u = \ln x$，将复合积分化为最基础的 $p$-反常积分进行分类讨论


\begin{solbox}[分步推导与答题规范]
\textbf{[1]} 凑微分与换元：


\begin{equation*}
\int_2^{+\infty} \frac{1}{x \ln^p x}\,\mathrm{d}x = \int_2^{+\infty} (\ln x)^{-p}\,\mathrm{d}(\ln x)
\end{equation*}


令 $u = \ln x$。当 $x \to 2$ 时 $u \to \ln 2 > 0$；当 $x \to +\infty$ 时 $u \to +\infty$：


\begin{equation*}
= \int_{\ln 2}^{+\infty} \frac{1}{u^p}\,\mathrm{d}u
\end{equation*}


\textbf{[2]} 分类讨论幂次 $p$：

当 $p > 1$ 时：


\begin{equation*}
\int_{\ln 2}^{+\infty} u^{-p}\,\mathrm{d}u = \left. \frac{u^{1-p}}{1-p} \right\vert{}_{\ln 2}^{+\infty} = 0 - \frac{(\ln 2)^{1-p}}{1-p} = \frac{(\ln 2)^{1-p}}{p-1} \quad (\text{收敛})
\end{equation*}


当 $p = 1$ 时：


\begin{equation*}
\int_{\ln 2}^{+\infty} \frac{1}{u}\,\mathrm{d}u = \left. \ln u \right\vert{}_{\ln 2}^{+\infty} = +\infty \quad (\text{发散})
\end{equation*}


当 $p < 1$ 时：


\begin{equation*}
\left. \frac{u^{1-p}}{1-p} \right\vert{}_{\ln 2}^{+\infty} = +\infty \quad (\text{发散})
\end{equation*}


\textbf{[3]} 结论：

当 $p > 1$ 时积分收敛；当 $p \leqslant 1$ 时积分发散。
\end{solbox}


\subsection{习题 8: 柯西-施瓦茨定积分不等式证明}

设函数$f(x)$在$[0, 1]$上连续，证明：


\begin{equation*}
\left(\int_0^1 f(x)\,\mathrm{d}x\right)^2 \le \int_0^1 f^2(x)\,\mathrm{d}x
\end{equation*}

并指出等号成立的充要条件。

\textbf{\color{DeepPurple}【思路点拨】} 构造实变量$t$的二次非负多项式积分$g(t) = \int_0^1 [t f(x) + 1]^2\mathrm{d}x \ge 0$，利用二次判别法$\Delta \le 0$证得不等式


\begin{solbox}[分步推导与答题规范]
\textbf{[1]} 构造辅助二次函数：

对任意实数$t$，考虑被积表达式$[t f(x) + 1]^2 \ge 0$。

由定积分保号性，积分值恒非负：


\begin{equation*}
g(t) = \int_0^1 [t f(x) + 1]^2\,\mathrm{d}x \ge 0
\end{equation*}


\textbf{[2]} 展开为关于$t$的一元二次多项式：


\begin{equation*}
g(t) = \int_0^1 [t^2 f^2(x) + 2t f(x) + 1]\,\mathrm{d}x = \left(\int_0^1 f^2(x)\,\mathrm{d}x\right) t^2 + 2\left(\int_0^1 f(x)\,\mathrm{d}x\right) t + \int_0^1 1\,\mathrm{d}x
\end{equation*}


记$A = \int_0^1 f^2(x)\,\mathrm{d}x$，$B = \int_0^1 f(x)\,\mathrm{d}x$，$C = 1$。

则$g(t) = A t^2 + 2B t + 1 \ge 0$对一切实数$t$均成立。

\textbf{[3]} 利用一元二次方程判别法：

二次三项式恒非负，其判别式必满足$\Delta \le 0$：


\begin{equation*}
\Delta = (2B)^2 - 4A \cdot 1 = 4B^2 - 4A \le 0 \implies B^2 \le A
\end{equation*}


即：


\begin{equation*}
\left(\int_0^1 f(x)\,\mathrm{d}x\right)^2 \le \int_0^1 f^2(x)\,\mathrm{d}x
\end{equation*}


\textbf{[4]} 等号成立条件：

等号成立当且仅当$\Delta = 0$，即存在实数$t_0$使得被积函数恒等于0：$t_0 f(x) + 1 \equiv 0 \implies f(x) \equiv -\frac{1}{t_0} = \text{常数}$。

故等号成立当且仅当$f(x)$为常数函数。证毕。
\end{solbox}


\subsection{习题 9: 含参变量变上限积分的高阶求导运算}

设$F(x) = \int_0^x (x - t)\sin t\,\mathrm{d}t$，求$F'(x)$与$F''(x)$。

\textbf{\color{DeepPurple}【思路点拨】} 积分内含有变上限自变量$x$，必须先将线性参变量$x$拆分提至积分号外，再应用乘积求导法则与微积分基本定理求一阶与二阶导数


\begin{solbox}[分步推导与答题规范]
\textbf{[1]} 参变量线性拆分：


\begin{equation*}
F(x) = \int_0^x (x\sin t - t\sin t)\,\mathrm{d}t = x \int_0^x \sin t\,\mathrm{d}t - \int_0^x t\sin t\,\mathrm{d}t
\end{equation*}


\textbf{[2]} 求一阶导数$F'(x)$：

对两项分别求导(第一项使用乘积求导法则)：


\begin{equation*}
F'(x) = \left[ 1 \cdot \int_0^x \sin t\,\mathrm{d}t + x \cdot \sin x \right] - x\sin x = \int_0^x \sin t\,\mathrm{d}t
\end{equation*}


\textbf{[3]} 求二阶导数$F''(x)$：

对$F'(x)$两边再求导：


\begin{equation*}
F''(x) = \frac{\mathrm{d}}{\mathrm{d}x} \int_0^x \sin t\,\mathrm{d}t = \sin x
\end{equation*}


\textbf{[4]} 故所求导数为：


\begin{equation*}
F'(x) = 1 - \cos x, \quad F''(x) = \sin x
\end{equation*}

\end{solbox}


\subsection{习题 10: 振荡型反常积分的狄利克雷审敛判定}

讨论反常积分$\int_1^{+\infty} \frac{\sin x}{x^p}\,\mathrm{d}x$在$p > 0$时的绝对收敛性与条件收敛性。

\textbf{\color{DeepPurple}【思路点拨】} 当$p>1$时利用绝对值比较法判定绝对收敛；当$0<p\le 1$时利用狄利克雷判别法证明收敛，并利用$\sin^2 x = \frac{1-\cos 2x}{2}$证明绝对值发散，从而判定为条件收敛


\begin{solbox}[分步推导与答题规范]
\textbf{[1]} 分析$p > 1$时的绝对收敛性：

考察绝对值积分$\int_1^{+\infty} \left\vert{}\frac{\sin x}{x^p}\right\vert{}\,\mathrm{d}x$。

因$\left\vert{}\frac{\sin x}{x^p}\right\vert{} \le \frac{1}{x^p}$，且$\int_1^{+\infty} \frac{1}{x^p}\,\mathrm{d}x$在$p > 1$时收敛。

由比较审敛法，$\int_1^{+\infty} \frac{\sin x}{x^p}\,\mathrm{d}x$在$p > 1$时\textbf{绝对收敛}。

\textbf{[2]} 分析$0 < p \le 1$时的收敛性(狄利克雷判别法)：

- 原函数有界性：$\left\vert{}\int_1^x \sin t\,\mathrm{d}t\right\vert{} = \vert{}-\cos x + \cos 1\vert{} \le 2$在$[1, +\infty)$上有界；
- 单调趋于0：当$p > 0$时，$g(x) = \frac{1}{x^p}$在$[1, +\infty)$上单调递减且$\lim_{x \to +\infty} \frac{1}{x^p} = 0$。

由狄利克雷审敛法，反常积分$\int_1^{+\infty} \frac{\sin x}{x^p}\,\mathrm{d}x$收敛。

\textbf{[3]} 分析$0 < p \le 1$时绝对值积分的发散性：


\begin{equation*}
\vert{}\sin x\vert{} \ge \sin^2 x = \frac{1 - \cos 2x}{2}
\end{equation*}



\begin{equation*}
\int_1^{+\infty} \frac{\vert{}\sin x\vert{}}{x^p}\,\mathrm{d}x \ge \frac{1}{2}\int_1^{+\infty} \frac{1}{x^p}\,\mathrm{d}x - \frac{1}{2}\int_1^{+\infty} \frac{\cos 2x}{x^p}\,\mathrm{d}x
\end{equation*}


由狄利克雷判别法，后半部分$\int_1^{+\infty} \frac{\cos 2x}{x^p}\,\mathrm{d}x$收敛；而前半部分$\int_1^{+\infty} \frac{1}{x^p}\,\mathrm{d}x$在$0 < p \le 1$时发散。

故绝对值积分发散。

\textbf{[4]} 结论：

当$p > 1$时，反常积分\textbf{绝对收敛}；当$0 < p \le 1$时，反常积分\textbf{条件收敛}。
\end{solbox}

\end{document}
